\documentclass[]{article}

\usepackage{amsmath}
\usepackage{multirow}
\usepackage{natbib}
\usepackage{graphicx} 


\begin{document}

Characterizing energy dissipation is fundamental to understanding the dynamics of physical systems, from microscopic resonators to large-scale civil structures. Experimentally, the total energy trajectory of a system provides direct insight into dissipative forces and allows for the estimation of key physical parameters, such as damping coefficients. However, such measurements are invariably corrupted by noise, which can obscure the true physical behavior. A common artifact of measurement noise is the emergence of an artificial, non-zero energy floor, particularly apparent in the late-time evolution of a system when its physical energy has decayed. This energy plateau arises because the total energy is typically calculated from squared quantities, such as position and velocity. Consequently, the variance of the zero-mean measurement noise manifests as a positive, constant offset in the calculated energy. This systematic bias corrupts the observed decay curve, making it difficult to accurately model the dissipation process and leading to erroneous estimates of physical parameters, especially in systems with low signal-to-noise ratios.

This paper presents an analytical deconvolution framework designed to isolate and remove this noise-induced energy bias from the measured energy decay of damped harmonic oscillators. The core of our approach is to characterize the noise by analyzing the system's behavior in the late-time regime, where the physical signal is negligible and the measurements are dominated by stationary noise. By calculating the empirical variances of the position ($\sigma_x^2$) and velocity ($\sigma_v^2$) signals in this phase, we can determine the constant energy offset introduced by the noise. For a harmonic oscillator with mass $m$ and spring constant $k$, this bias term is given by $\Delta E_{\text{noise}} = \frac{1}{2}(k\sigma_x^2 + m\sigma_v^2)$.

We demonstrate that by subtracting this analytically derived bias from the total measured energy, we obtain a corrected energy trajectory that faithfully represents the true physical dissipation. This corrected trajectory follows the expected exponential decay law, $E(t) = E_0 \exp(-2\gamma t)$, down to zero, even when the original data exhibited a significant energy floor. We validate our method by applying it to a dataset of simulated damped oscillators and showing that a non-linear least-squares fit of the exponential model to the corrected data allows for the recovery of the theoretical damping rate $\gamma$ with high precision. This framework provides a robust and straightforward method for distinguishing measurement artifacts from physical reality, thereby enabling more accurate parameter estimation and a clearer validation of energy dissipation models from noisy experimental data.

\end{document}
                